\documentclass[12pt, a4paper]{article}

\usepackage[utf8]{inputenc}
\usepackage{amsmath, amssymb}
\usepackage{graphicx}
\usepackage{booktabs}
\usepackage{geometry}
\usepackage{hyperref}
\usepackage{float}

\geometry{margin=1in}

\title{\textbf{Model Predictive Control Architectures for Topologically Coupled Agro-Systems}}
\author{} % Add your name or Tara's name here
\date{\today}

\begin{document}

\maketitle

\begin{abstract}
This report outlines four advanced Model Predictive Control (MPC) architectures designed to optimize irrigation in a topographically complex, 130-agent agricultural environment. Because accurate surface hydrology requires a daily cascading topological sort, the agents are strictly coupled within each time step. A standard 14-day prediction horizon generates over 1,820 interdependent decision variables, causing generic dense solvers to scale cubically $\mathcal{O}(n^3)$ and fail. The architectures proposed herein mitigate this computational bottleneck through sparsity exploitation, differentiable simulation, hierarchical clustering, and neural imitation, providing a clear roadmap for scalable closed-loop control.
\end{abstract}

\section{Introduction to Model Predictive Control}
\label{sec:intro_mpc}

Model Predictive Control (MPC) is an advanced feedback control strategy that leverages an explicit mathematical model of a system to predict its future evolution and optimize control decisions. Unlike traditional reactive controllers that respond only to current errors, MPC is inherently proactive. By incorporating future trajectories and measurable disturbances—such as 14-day weather forecasts—MPC anticipates system behavior and calculates a sequence of optimal control actions over a defined prediction horizon ($H_p$).



The operational mechanism of MPC relies on the ``receding horizon'' principle, which executes via the following loop:
\begin{enumerate}
    \item \textbf{Measurement:} At the current discrete time step $k$, the controller receives the measured or estimated state of the system ($x(k)$).
    \item \textbf{Optimization:} The controller solves a constrained optimization problem to find an optimal sequence of future control actions $\{u(k), u(k+1), \dots, u(k+H_p-1)\}$ that minimizes a specified cost function while strictly obeying system constraints.
    \item \textbf{Execution:} Only the first calculated control action, $u(k)$, is applied to the physical system. 
    \item \textbf{Recession:} At the next time step, $k+1$, the horizon shifts forward by one step, new measurements are taken, and the entire optimization process is repeated. 
\end{enumerate}

This continuous closed-loop feedback allows the MPC to seamlessly correct for unmodeled dynamics, prediction errors, and unexpected environmental disturbances. 

To function effectively, an MPC requires three rigorously defined components:
\begin{itemize}
    \item \textbf{The Predictive Model:} A discrete-time dynamical representation of the system (e.g., the Agent-Based crop-soil hydrology model).
    \item \textbf{The Objective Function:} A mathematical cost to be minimized, which balances the desired performance (maximizing final crop biomass) against the control effort (minimizing total water consumption).
    \item \textbf{Constraints:} Strict physical or operational limits, such as maximum soil saturation capacities, maximum daily irrigation bounds, and global volumetric water budgets.
\end{itemize}

While the fundamental theory of MPC is mathematically elegant, its practical implementation becomes highly complex as the physical scale of the system grows. In networked cyber-physical systems—such as a 130-agent topographical farm where water routes continuously between spatial nodes—a standard centralized MPC requires solving thousands of interdependent decision variables ($H_p \times N$) at every single time step. 

To achieve real-time control without exceeding the computational limits of standard hardware, the fundamental MPC algorithm must be adapted into specialized structural architectures. The subsequent sections detail how different architectures handle this spatial and computational complexity to maintain optimal irrigation policies.

\section{The Computational Control Challenge}
\label{sec:challenge}
The agricultural environment (Gilan province model) simulates 130 agents interacting over a directed elevation graph. To accurately simulate overland flow using the SCS Curve Number method, the system utilizes a single-day cascade routing logic. This establishes a strict intra-day dependency: a decision to irrigate an uphill agent instantaneously impacts the available surface water ($W_{\text{surf}}$) and waterlogging stress ($h_6$) of downhill sink agents.

When formulating the MPC over a prediction horizon ($H_p = 14$), the resulting non-linear programming (NLP) problem contains $H_p \times N = 1,820$ decision variables. Traditional sequential quadratic programming (SQP) algorithms evaluate dense Jacobians, causing memory and execution time to scale cubically. To achieve real-time control, the architecture must bypass dense matrix operations while strictly respecting the global water budget and the cascade physics.

\section{Architecture 1: Centralized Sparse NLP}
\label{sec:casadi}
This architecture retains the centralized mathematical formulation—solving the entire 130-agent grid simultaneously—but leverages Advanced Algorithmic Differentiation (AD) and sparse matrix solvers to handle the computational load.

\begin{itemize}
    \item \textbf{Methodology:} The discrete-time physical equations are translated into symbolic variables using the \textbf{CasADi} framework. Because water only flows to direct topographical neighbors, the system's Jacobian is highly sparse. The problem is solved using the \textbf{IPOPT} interior-point solver, which natively exploits this sparsity.
    \item \textbf{Advantages:} It perfectly respects the cascade physics and provides the absolute mathematical ``ground truth'' for the optimal irrigation schedule. Constraints, such as a strict volumetric water budget, are easily enforced as linear boundaries.
    \item \textbf{Disadvantages:} Formulating the symbolic equations requires meticulous engineering. Even with sparse linear algebra, the computational overhead may remain too high for deployment on low-power edge devices.
\end{itemize}

\section{Architecture 2: Differentiable Simulation}
\label{sec:jax}
Rather than constructing a traditional mathematical optimization problem, this architecture frames the MPC as a computational graph, treating the physical simulation identically to a deep learning layer.

\begin{itemize}
    \item \textbf{Methodology:} The agent-based model is ported into a differentiable tensor library such as \textbf{JAX} or \textbf{PyTorch}. The 14-day prediction horizon is unrolled in time, and the final crop yield acts as the loss function. The daily irrigation array ($u$) is treated as a set of network weights. Optimization is performed via Gradient Descent (e.g., Adam) by backpropagating through the physical equations.
    \item \textbf{Advantages:} It bypasses manual Jacobian derivation and executes exceptionally fast on GPU hardware via vectorization. It seamlessly bridges classical control with modern AI infrastructure.
    \item \textbf{Disadvantages:} Gradient descent struggles with ``hard'' constraints (like an absolute water limit). Constraints must be enforced via heavy penalty terms in the objective function, which requires careful tuning.
\end{itemize}

\section{Architecture 3: Hierarchical MPC (HMPC)}
\label{sec:hmpc}
To handle massive physical scale, this architecture splits the optimization into two distinct spatial and temporal layers, mirroring the physical operation of industrial manifold valves.

\begin{itemize}
    \item \textbf{Methodology:} The 130 agents are clustered into Topographical Management Zones (e.g., Highlands, Mid-Slopes, Sinks). 
    \begin{enumerate}
        \item \textit{Macro-Level:} A high-level MPC allocates the total daily water budget among the zones over the 14-day horizon.
        \item \textit{Micro-Level:} A localized algorithm distributes the zone's daily budget to individual agents based on real-time soil moisture deficits ($x_1$).
    \end{enumerate}
    \item \textbf{Advantages:} Exponentially reduces the decision space for the optimizer (e.g., from 130 variables to 3 variables per time step). Highly scalable and realistic for field deployment.
    \item \textbf{Disadvantages:} Inherently suboptimal compared to a centralized solver. The macro-level relies on aggregated zone data, potentially obscuring localized drought or waterlogging pockets.
\end{itemize}

\section{Architecture 4: Imitation / Neural MPC}
\label{sec:imitation}
This architecture leverages deep learning to distill the complex optimal control policy into a lightweight, real-time function approximator.

\begin{itemize}
    \item \textbf{Methodology:} A heavy, centralized MPC (Architecture 1) is executed offline thousands of times across synthesized historical weather scenarios to generate optimal state-action pairs. A Deep Neural Network (DNN) is then trained via Supervised Learning to map the current field state and weather forecast directly to the optimal irrigation commands.
    \item \textbf{Advantages:} Eliminates runtime optimization entirely. The deployed neural network executes the control policy in milliseconds, making it ideal for edge computing hardware (e.g., Raspberry Pi).
    \item \textbf{Disadvantages:} The controller's reliability is strictly bounded by the diversity of the offline training data. Out-of-distribution weather events may cause unpredictable irrigation behavior.
\end{itemize}

\section{Evaluation and Comparison Framework}
\label{sec:evaluation}
To scientifically benchmark the selected architectures, they must be evaluated against a unified set of criteria over canonical historical climate scenarios (e.g., Wet 2024, Dry 2022). The evaluation metrics are defined in Table \ref{tab:metrics}.

\begin{table}[H]
\centering
\caption{Proposed Controller Evaluation Metrics}
\label{tab:metrics}
\begin{tabular}{@{}llp{8cm}@{}}
\toprule
\textbf{Category} & \textbf{Metric} & \textbf{Description} \\ \midrule
Agronomic & Mean Yield ($Y$) & Average final biomass converted to kg/ha. \\
Agronomic & WUE & Water Use Efficiency (biomass produced per mm applied). \\
Agronomic & Spatial Equity (CV) & Coefficient of Variation of final biomass across the 130 agents. \\
Computational & Mean Solve Time ($\tau_{\text{mean}}$) & Average milliseconds required to compute a daily action step. \\
Computational & Max Latency ($\tau_{\text{max}}$) & Worst-case computational time (critical for edge feasibility). \\
Robustness & Noise Degradation & Yield loss when subjected to a stochastically noisy weather forecast. \\ \bottomrule
\end{tabular}
\end{table}

\section{Conclusion}
\label{sec:conclusion}
The topological coupling of the Gilan farm model renders standard SQP optimization obsolete. To establish a rigorous mathematical baseline, the Centralized Sparse NLP (Architecture 1) should be implemented first. Following this, exploring Differentiable Simulation (Architecture 2) or Neural Imitation (Architecture 4) provides a highly novel, scalable solution that directly sets the stage for advanced Reinforcement Learning deployments.

\end{document}